\section{What Flux is}
\label{sec:short:what}

Concrete before abstract. Everything in this section was measured on 2026-09-03 from the network's
own public interfaces, and can be measured again by anyone who runs the scripts in the repository.

\subsection{A cloud made of collateralised machines}

\Flux{} is a decentralized cloud: a network of independently operated machines that run other
people's applications, coordinated by a blockchain rather than by a company. On the day of
measurement it consisted of \textbf{\num{6489} nodes} running \textbf{\num{1195} deployed
applications}.

What makes a machine a node is not an agreement with anyone. It is a lock. An operator locks a
fixed amount of \flux{} as collateral, runs the node software, and from then on the chain itself
decides whether that machine is eligible to host and to be paid. Leave, and the collateral unlocks.
Misbehave in ways the chain can see, and it does not.

Three tiers, three collateral sizes, three hardware floors:

\begin{center}\small
\begin{tabular}{@{}lrrr@{}}
\toprule
Tier & Collateral & Nodes today & Share of voting weight \\
\midrule
Cumulus & \num{1000}~\flux & \num{3247} & \SI{3.67}{\percent} \\
Nimbus & \num{12500}~\flux & \num{1615} & \SI{22.81}{\percent} \\
Stratus & \num{40000}~\flux & \num{1627} & \SI{73.52}{\percent} \\
\midrule
\multicolumn{2}{@{}l}{Total locked} & \num{6489} & \num{88514500}~\flux \\
\bottomrule
\end{tabular}
\end{center}

Two things in that table are worth noticing before moving on. First, the collateral is real and
large: \num{88514500}~\flux{} is locked across the fleet, about a fifth of everything ever issued.
Second, weight follows collateral, not headcount --- half the nodes are Cumulus, but they carry under
four percent of the weight. That fact returns in \Cref{sec:short:limits}, where it stops being a
curiosity.

\subsection{A chain that pays for capacity, not for hashing}

The network runs on its own chain. Until October 2025 that chain was proof-of-work; at block
\num{2020000} it switched to \emph{Proof of Node}, in which the right to produce a block is drawn
by lot from the collateralised nodes rather than won by burning electricity. Blocks arrive every
thirty seconds. Node income is a position in a rotation, not a lottery ticket.

The consequence for a reader is simple: the thing the chain rewards is \emph{being available with
collateral at stake}, which is exactly the thing a cloud needs its machines to be.

\subsection{What actually runs on it}

Applications are containers, described by a specification the owner publishes to the chain and pays
for by the block. Three-quarters of deployed applications (\num{873} of \num{1195}) are on the
current specification version, and the mix is more general than a newcomer might guess: ordinary web
services and content sites; blockchain infrastructure, including full nodes for other chains;
and a substantial block of scientific-computing work --- \num{27} protein-folding and
structural-biology deployments that between them hold roughly \SI{9}{\percent} of the fleet's
committed CPU.

That last point matters enough to come back to in \Cref{sec:short:whatfor}. For now: this is a
running, general-purpose cloud, not a proposal for one.

\limitnote{Capacity is measurable; utilisation is not. The interfaces this section was measured from
report what each node \emph{has} and what is \emph{deployed}, not how busy anything is. The paper
states the counts it can stand behind and does not estimate the number it cannot.}
